\scanpage{78}
% archive sheet number: 79
est de rang fini \inserted{(vérifier d'après cas décomposable)} (voir Chap I, \S~1, N\textsuperscript{o} 2, définition 1), et reste donc vraie par \unsure{raison|raisons} de continuité dans le cas général. Conjointement avec (13), on trouve

\uline{Corollaire 2} Sous les conditions précédentes, on a

\[
|\alpha_n(u_1,\ldots u_n)|\leq\frac1{n!}\lVert u_1\rVert\ldots\lVert u_n\rVert\tag*{(15)}
\]

(Cela résulte aussi du fait, déjà \unsure{signalé}, que le coëfficient \unsure{$k_n$|$h_n$} du Chap I, \S~1, N\textsuperscript{o} 2, prop 1, est ici égal à 1).

\uline{Corollaire 3} Soit \deleted{$u\in L^1(E)$, où $E$ est un espace de Hilb} \inserted{\unsure{$u$ un} op. de Fredholm dans un} espace de Hilbert $E$. \deleted{Alors,} Si $u$ est un op. de Fredholm, il en est de même de $\overset{n}{\Lambda}u$, et on a

\[
\lVert\overset{n}{\Lambda}u\rVert_1=\alpha_n\deleted{\unsure{(\sqrt{u^*u})}}(|u|)=\alpha_n(\sqrt{\unsure{u^*u|u^*}})
\tag*{(16)}
\]

Si $u$ est un op. de H--S., il en est de même de $\overset{n}{\Lambda}u$, et on a

\[
\lVert\overset{n}{\Lambda}u\rVert_{\unsure{2|1}}=\sqrt{\alpha_n(u^*u)}\tag*{(\unsure{17|14})}
\]

On a, si $u\in L^1(E)$, $\lVert\overset{n}{\Lambda}u\rVert_1=\operatorname{Tr}\deleted{\unsure{\sqrt{\vphantom{(\overset{n}{\Lambda}u)^*}}}}\bigl((\overset{n}{\Lambda}u)^*(\overset{n}{\Lambda}u)\bigr)^{1/2}$ (N\textsuperscript{o} 1, formule (22)), or, (prop 1) on a $\deleted{\unsure{\sqrt{\vphantom{(\overset{n}{\Lambda}u)^*}}}}\bigl((\overset{n}{\Lambda}u)^*(\overset{n}{\Lambda}u)\bigr)^{1/2}=\overset{n}{\Lambda}\bigl((u^*u)^{1/2}\bigr)$, d'où \deleted{(16) en vertu de (17)} \inserted{$\lVert\overset{n}{\Lambda}u\rVert_1=\operatorname{Tr}\overset{n}{\Lambda}((u^*u)^{1/2})=\alpha_n(\sqrt{u^*u})$} (en vertu de (14)).

\deleted{Si $u\in L^2(E)$, \inserted{alors} on a encore $(\overset{n}{\Lambda}u)^*(\overset{n}{\Lambda}u)=\overset{n}{\Lambda}(u^*u)$ \unsure{où on} \unsure{avait} $u^*u\in L^1(E)$ \inserted{est op. de Fredholm} (d'où $(\overset{n}{\Lambda}u)^*(\overset{n}{\Lambda}u)=\overset{n}{\Lambda}(u^*u)$ est $\in L^1(E)$, et $\overset{n}{\Lambda}u$ \unsure{par suite}}

Si $u$ est un op. de H.--S., alors $u^*u$ est un op. de Fredholm (N\textsuperscript{o} 1, th.~4), donc $(\overset{n}{\Lambda}u)^*(\overset{n}{\Lambda}u)=\overset{n}{\Lambda}(u^*u)$ est un op. de Fredholm, \deleted{et} donc $\overset{n}{\Lambda}u$ est un op. de H.--S. On a (N\textsuperscript{o} 1, th.~3, corollaire 3)

\[
(\lVert\overset{n}{\Lambda}u\rVert_2)^2=\operatorname{Tr}(\overset{n}{\Lambda}u)^*(\overset{n}{\Lambda}u)=\alpha_n(u^*u)\qquad\text{(formule (14)) d'où (17)}.
\]
